\documentclass[11pt]{article}
\usepackage{amsmath,amssymb,amsthm}
\newtheorem{proposition}{Proposition}
\newcommand{\E}{\mathbb E}
\title{A vacuum-relative local-frequency inequality for the original Wilson transfer}\author{}\date{}
\begin{document}\maketitle
\paragraph{Exact physical operator and conventions.}
Use the ORIGINAL $SU(N)$ Wilson transfer on a periodic spatial lattice with sides
at least three, $A=M Q P M$, $M=e^{-V/2}$, $V=B\sum_{p\,\rm spatial}s(U_p)$,
$s(U)=1-N^{-1}\operatorname{Re}\operatorname{Tr}U$. Temporal plaquettes have the
same normalized coefficient $B$. The site-gauge projection $P$ is retained.
Let $\lambda_0>0$, $\psi>0$, $\|\psi\|_{L^2(dX)}=1$ be the actual leading
eigenpair, and put $\pi=\psi^2dX$. The ORIGINAL ground-state Markov operator is
\[
 (\mathcal T f)(X)=\frac1{\lambda_0\psi(X)}
          \int A(X,Y)\psi(Y)f(Y)dY.                    \tag{1}
\]
It is positive self-adjoint on $L^2(\pi)$ and preserves the physical subspace.
Its positive physical eigenvalues are $r_j=\lambda_j(A)/\lambda_0$.
All estimates below are on that space. No temporal scalar or unknown vacuum
energy is substituted for $\lambda_0$.

Choose anti-Hermitian generators with $-\operatorname{Tr}T^AT^B=\delta^{AB}$.
For one spatial link $b$, let $\nabla_b=(X_b^A)_A$, $\Delta_b=\sum_A(X_b^A)^2$,
and $C_F=(N^2-1)/N$. Define the nonnegative ground-weighted one-link operator
$D_b$ by its closed form $\mathcal E_b(f)=\int|\nabla_b f|^2d\pi$.
Conditioned on all other spatial links it is the weighted Laplacian on the
actual compact one-link fiber, with density proportional to $\psi^2$.
This definition is not a free-link Wilson conditional or a guessed native law.

\begin{proposition}[Exact vacuum-relative gradient smoothing]
For every smooth physical $f$,
\[
 \mathcal E_b(\mathcal T f)
 \le K_B\big(\|f\|_\pi^2-\|\mathcal T f\|_\pi^2\big),
 \qquad K_B=B^2/N.                                    \tag{2}
\]
The constant is independent of spatial volume, cutoff and the actual value of
$\lambda_0$.
\end{proposition}
\begin{proof}
Write $q_P(X,Y)$ for the kernel of $QP$. It is the original temporal-link gauge
average of product Wilson temporal weights. Differentiating in $X_b$ touches
one temporal plaquette for each gauge configuration. Its score vector obeys
$|\nabla_b\log k|\le B/\sqrt N$: the traceless anti-Hermitian projection of a
unitary matrix has Hilbert--Schmidt norm at most $\sqrt N$.
A positive gauge average preserves this bound, so
$h_b=\nabla_b\log q_P$ satisfies $|h_b|\le B/\sqrt N$.

Differentiate the eigenfunction normalization in (1). The deterministic
$\nabla_b\log M(X)$ cancels, as does the normalizing score from $\psi(X)$:
\[
 \nabla_b\mathcal T f
 =\operatorname{Cov}_{p_0(X,\cdot)}(f,h_b),\qquad
 p_0(X,Y)=A(X,Y)\psi(Y)/(\lambda_0\psi(X)).             \tag{3}
\]
Vector Cauchy--Schwarz gives
$|\nabla_b\mathcal T f|^2\le K_B\{\mathcal T|f|^2-|\mathcal T f|^2\}$.
Integrate using the ACTUAL stationarity $\pi\mathcal T=\pi$.
For complex $f$ use the usual absolute-value conditional variance.
No tensorization or mixing assumption enters.
\end{proof}

\paragraph{An actual conditional-frequency projection bound.}
Let $E_b(\Lambda)=\boldsymbol1_{(\Lambda,\infty)}(D_b)$, $\Lambda>0$.
It is a direct-integral conditional spectral projection; gauge covariance of
$D_b$ preserves the physical space. From $D_b\ge\Lambda E_b(\Lambda)$ and (2),
\[
 \boxed{\quad
 \mathcal T E_b(\Lambda)\mathcal T
       \le\frac{B^2}{N\Lambda}(I-\mathcal T^2).
 \quad}                                               \tag{4}
\]
This is a ground-relative operator inequality, with the vacuum canceled BEFORE
any trace bound. For the spectral subspace $S_E$ of true physical energies
$-a^{-1}\log r_j\le E$, invert $\mathcal T$ on this finite-dimensional subspace
to obtain
\[
 S_E E_b(\Lambda)S_E
 \le\frac{B^2}{N\Lambda}(e^{2aE}-1)S_E.               \tag{5}
\]
Summing the left side of (5) over links costs its literal $L$. For fixed $E$
and $L\asymp\operatorname{Vol}\Sigma/a^3$, this gives at most
$C\operatorname{Vol}\Sigma B^2E/(\Lambda a^2)$. With $\Lambda\asymp B$
the single-link $O(aBE)$ error is summable over geometric cutoffs, but its
all-link sum is not. Taking $\Lambda\gtrsim B^2/a^2$ controls that sum only by
allowing a growing local representation cutoff; it does not give a fixed physical
coarse state count.

In particular the ground vector $f=1$ has zero high-frequency weight. For a
fixed physical excitation energy and $\Lambda\asymp B$, the right side is
$O(aBE)$, uniformly at each link. This is about ground-adapted conditional angular
frequency, not an assertion about spatial Fourier modes or free torons.

A trace consequence retaining all physical states is, for integer $m\ge3$,
\[
 \operatorname{Tr}\{E_b(\Lambda)\mathcal T^m\}
 \le\frac{B^2}{N\Lambda}
   [Q_{\rm occ}((m-2)a)-Q_{\rm occ}(ma)].              \tag{6}
\]
Multiply (4) by the positive $\mathcal T^{m-2}$ and take its trace; cyclicity
justifies the left side. The restriction $m\ge3$ keeps both traces finite.
Equation (6) is not a uniform upper bound for $Q_{\rm occ}$ by itself.

\paragraph{The local conditional frequency has a controlled Casimir count.}
There is a pointwise lower Hessian-trace bound with only a LOCAL cost:
\[
 \Delta_b\log\psi\ge-3 B C_F.                         \tag{7}
\]
To verify it, use $\psi=M\lambda_0^{-1}\int q_P(X,Y)M(Y)\psi(Y)dY$.
For every temporal gauge configuration the log temporal weight has
$\Delta_b\log k\ge-B C_F$. Passing through a positive integral adds the
nonnegative trace of the score covariance. Each of the four spatial plaquettes
incident to $b$ contributes at worst $-(B/2)C_F$ to $\Delta_b\log M$.
Their total is $-2BC_F$, giving (7). All counts are original three-dimensional
spatial counts; no microscopic boundary or gauge orbit count was omitted.

On a fixed conditional fiber, multiplication by its normalized $\psi$ is
unitary from the weighted space to Haar space and gives
\[
 D_b\ \simeq\ -\Delta_b+\frac{\Delta_b\psi}{\psi}
 =-\Delta_b+\Delta_b\log\psi+|\nabla_b\log\psi|^2
 \ge-\Delta_b-3BC_F.                                 \tag{8}
\]
Thus the number of local conditional eigenvalues below $\Lambda$ is at most
\[
 \sum_{R:C_R\le\Lambda+3BC_F}d_R^2,                  \tag{9}
\]
uniformly in the actual bath. For $SU(2)$, $C_d=(d^2-1)/2$, so this is
$\sum_{d\le\sqrt{1+2\Lambda+9B}}d^2$.
The original full physical eigenfunction $\phi=\psi f$ with energy at most $E$
also obeys the genuine Haar-link kinetic estimate
\[
 \langle\phi,(-\Delta_b)\phi\rangle
 \le[3BC_F+(B^2/N)(e^{2aE}-1)]\|\phi\|^2.             \tag{10}
\]
In contrast to (5), (10) includes the legitimate ground-state kinetic baseline.
No inverse ground overlap or global vacuum energy enters either estimate.

\paragraph{Global counting attempt and its remaining cutoff cost.}
Sum (2) over the $L$ spatial links. The global weighted operator
$D=\sum_bD_b$ is conjugate to
$-\sum_b\Delta_b+(\sum_b\Delta_b\psi)/\psi$
and is bounded below by the physical Haar Casimir minus $3BC_FL$.
Min--max on $S_E$ therefore gives the actual interacting ground-relative count
\[
 N_A(E)\le N_{\rm Haar,phys}
 \left(L[3BC_F+(B^2/N)(e^{2aE}-1)]\right).             \tag{11}
\]
This is proved without a ground-energy approximation. It is still insufficient
for a fixed-physical-volume occupation estimate. For example, the explicit
$SU(2)$ product Peter--Weyl count gives the upper estimate
\[
 \log N_A(E)\le L\{C+\tfrac32\log B
                      +\tfrac B2(e^{2aE}-1)\},\qquad B\ge1,    \tag{12}
\]
by exponential counting at parameter $1/B$ and
$\sum_d d^2e^{-(d^2-1)/(2B)}\le C B^{3/2}$.
The full product space bounds the physical count from above; it does not replace
Gauss projection with equality. At fixed $E$ the $L\log B$ loss remains cutoff
extensive. Bounds (11)--(12) do NOT assert that this many low-energy states
actually exist.

The conditional low-frequency projectors at different links depend on the
entire bath and need not commute. Their finite fiber ranks in (9) cannot be
multiplied to infer a physical coarse-mode tensor product. Moreover a linkwise
angular-frequency estimate does not control rapid spatial oscillations of
collective low angular-frequency coordinates. The unresolved spatial stiffness
and conditional-basis compatibility are the precise missing improvements.

\paragraph{A stronger averaged Ward identity, and why it does not yet tensorize (2).}
There is an additional exact improvement for deterministic multi-link tangent
fields $w$. Write $X_w=\sum_{b,A}w_b^A X_b^A$ and
$h_w=X_w\log q_P$. Differentiating the same ground equation twice and integrating
by parts in ORIGINAL Haar measure yields
\[
 \E_\pi\operatorname{Var}_{p_0(X,\cdot)}h_w
       +2\|X_w\log\psi\|_\pi^2
 =-\E_\pi X_w^2\log M
       -\E_{\pi p_0}X_w^2\log q_P
 \le\frac{9B}{N}\sum_b|w_b|^2.                        \tag{13}
\]
For a tangent supported at one link the constant improves from $9$ to $3$.
Indeed $|X_w^2s(U_p)|\le N^{-1}(\sum_{b\in p}|w_b|)^2$;
each spatial plaquette has four links and each link meets four plaquettes,
so $X_w^2\log M\ge-8B\|w\|^2/N$. Before gauge averaging the temporal
Hessian is at least $-B\|w\|^2/N$; positive gauge averaging adds nonnegative
score covariance. Finally $\E_\pi X_w^2\log\psi=-2\|X_w\log\psi\|_\pi^2$.
This proves (13) with no total-volume factor in its deterministic directional
constant and no guessed conditional measure.

Equation (13) controls the AVERAGED temporal Fisher matrix by $9B I/N$,
considerably better than a pointwise $B^2 L$ estimate. It cannot simply replace
the coefficient in the summed version of (2): the latter integrand contains
$\operatorname{Var}_{p_0}(f\mid X)$ multiplied by the pointwise Fisher matrix.
These quantities are correlated. In a spectral trace the multiplier becomes a
conditional return/occupation diagonal, precisely a quantity still needing
control. Replacing an average of their product by a product of averages would
be an unproved step. This locates an exact possible use of spatial/coarse
information without assuming that missing estimate.

\paragraph{Consistency checks and scope.}
For pure kinetic $SU(2)$ convolution, (2) yields only a rational spectral upper
bound of the form $q_d^2 C_d\le K_B(1-q_d^2)$; the exact factorial high-$d$
Wilson tail is stronger, so no false Gaussian-Casimir domination is introduced.
Actual constant/holonomy sectors have not been removed: only gauge-projected
physical states are counted, and (5) is zero on the true ground state without
claiming that all global excitations have a gap. The derivations require neither
unavailable R-I domains nor any native convex approximation.

Equations (2),(4),(5),(7)--(11) are genuine original-transfer estimates. They
advance local ground-relative high-frequency control, but the remaining
cutoff-extensive count (12) prevents an S94 closure. To improve it one needs
actual spatial/coarse compatibility or a stronger scale-resolved response,
not an assumed independent tensor product of conditional fibers.
\end{document}
